\begin{textsample}{NEG Dim 6 – global\_north\_2025\_03 – Score 1.00 – global\_north\_2025\_03/121873600.txt}  \label{ex:f6_neg_172}
Jewish lawmakers want less discussion of Israel and Palestine in soon-to-be-required ethnic studies courses at California public schools, which they say have occasionally veered into antisemitism.
Members of the Legislative Jewish Caucus and Bay Area Democrats introduced a bill last month that would require the state Board of Education to adopt standards for ethnic studies by January 2028.
A state educational commission would recommend course materials to the board.
The bill would require the Board of Education to monitor and post online the course materials being taught in classrooms throughout the state to ensure they don't reflect"bias, bigotry, or discrimination"and focus coursework on"the domestic experience and stories of historically marginalized peoples in American society,"not conflicts abroad."We don't think that ethnic studies is a foreign policy discipline,"said Tyler Gregory, a supporter of the bill who leads the Jewish Community Relations Council ' s Bay Area branch, which provided input to lawmakers on the plan."So to that end, yes, @ @ @ @ @ @ @ @ @ @ the appropriate place to teach it."That ' s sparking outrage among some teachers and Muslim-American advocacy groups."We can not teach an accurate picture of many communities ' lived experiences domestically if we are unable to discuss what is happening internationally,"said Zahra Billoo, executive director of the Bay Area branch of the Council on American-Islamic Relations, which opposes the bill and has gathered 1,100 signatures on a petition against it."For example, what is happening in Palestine right now is an integral part of the Arab-American experience."Legislation passed in 2021 requires California public high schools to offer a course in ethnic studies by the 2025-26 school year, and by 2030, students won't be able to graduate without it.
California teachers have more leeway to teach ethnic studies than other subjects.
A hard-fought model curriculum that was hammered out by experts and officials in 2021 is only a set of guidelines for teachers, according to the bill ' s sponsors.
The bill is the latest twist @ @ @ @ @ @ @ @ @ @, which has become a lightning rod for culture war arguments about diversity, history and racial identity -- especially during the Israel-Gaza war that erupted in October 2023.
Ethnic studies teachers aim to promote unity and shared understanding through discussions of race and ethnicity in the U.S., with a focus on African Americans, Latino and Chicanx people, Asian Americans and Pacific Islanders, and Native Americans, according to the model curriculum.
But some teachers have focused on the searing conflict between Israel and Palestine in discussions of colonialism and oppression.
The bill, AB 1468, was introduced by Assembly and Jewish Caucus members Rick Chavez Zbur of Los Angeles and Dawn Addis of San Luis Obispo, both Democrats who introduced a similar bill last year that later was withdrawn.
More than two dozen other Democrats in the state Assembly and Senate have signed onto the bill so far -- a significant show of support.
State Senator Josh Becker, a Menlo Park Democrat who backs the bill, said it wouldn't"ban @ @ @ @ @ @ @ @ @ @ be better settings to discuss Israel and Palestine, and educators would benefit from having clear expectations about what to teach -- and how."There ' s opportunity in World History to dig into all that discussion -- the complicated history of the Middle East,"he said.
Becker said one reason he supports the bill is a controversy over class materials in 2023 that some students and parents said was antisemitic and inaccurate at Menlo-Atherton High School, which Becker ' s son attended.
The divisive material included a drawing of a hand controlling a puppet with strings, which Becker and others consider an antisemitic trope portraying Jews as wielding outsized influence in the world.
The teacher, Chloe Gentile-Montgomery, no longer works there.
Becker said course standards might have protected her and other teachers."All of this stuff is unnecessary,"he said."If we had standards, hopefully we wouldn't be having these divisive fights."In a recent interview, Gentile-Montgomery said the class materials -- which @ @ @ @ @ @ @ @ @ @ context and that she never intended any harm, but she does regret the puppet imagery and"would never use that slide again."She said the situation left her rattled and nervous about teaching ethnic studies again and she supports content and teaching standards for ethnic studies."I ' m afraid of making missteps, I ' m afraid of using certain words,"she said."Our students who are seeing this every day on social media and have genuine questions about it,"she said."I feel like that also spreads the message of, ' just don't talk about it. '"It ' s unclear if the proposal will meet other opposition.
The California Legislative Black Caucus, whose members helped push a 2021 compromise on ethnic studies guidelines across the finish line, and the California Teachers Association haven't taken a position.
The CTA did oppose an attempt by legislators last year to strip"bias"from ethnic studies courses and require community review of course materials. @ @ @ @ @ @ @ @ @ @ find support for it, despite the backing of State Superintendent of Instruction Tony Thurmond and Gov.
Gavin Newsom.
However, college-level ethnic studies educators have quickly condemned the bill.
University of California faculty members from the Ethnic Studies Faculty Council deemed the proposal"chilling censorship"and criticized calls for more state oversight."What ethnic studies is really about is fighting racism,"said Keith Miyake, an ethnic studies professor at UC Riverside."It ' s about fighting discrimination and it ' s done in a critical way that bills like AB 1468 are designed to eliminate."

% matched lemmas: 
\end{textsample}
